\label{ch:spectral}

\section{Representaci\'on operatoria de los caracteres de Dirichlet}

Las constantes de este cap\'itulo se agrupan por compartir una construcci\'on
espectral independiente de sus expansiones decimales y libre de valores
objetivo. Sea

\[
Ne_n=ne_n,\qquad n\in\N,
\]

en \(\ell^2(\N)\). Un car\'acter residual \(\chi\) define el operador diagonal \(Q_\chi e_n=\chi(n)e_n\); siempre que la potencia sea traza--clase,

\begin{equation}
\Tr(N^{-s}Q_\chi)=L(s,\chi).
\label{eq:L-trace}
\end{equation}

La ecuaci\'on conserva el m\'odulo, la orientaci\'on y la procedencia de cada
car\'acter; por tanto, no identifica objetos residuales distintos.

\section{Construcci\'on HMT de los caracteres residuales}

El operador diagonal anterior no constituye el dato algebraico inicial. La
cadena de selecci\'on es

\[
\begin{aligned}
\APP&\longrightarrow\TRIT\longrightarrow\TPK
\longrightarrow\{\text{residuo m\'odulo }3,\ C_P\}\\
&\longrightarrow\{\chi_{-3},\ J=C_P\bmod3\}
\longrightarrow\{\chi_{-3},\ J^2=-I,\chi_{-4}\}\\
&\longrightarrow Q_3,Q_4.
\end{aligned}
\]

El funcional residual tri\'adico determina la diferencia entre las dos clases
orientadas no nulas. En el sector de incidencia se emplea la matriz de conferencia
sim\'etrica de Paley de orden seis,
\[
C_P^{\mathsf T}=C_P,\qquad C_PC_P^{\mathsf T}=5I_6.
\]
Su reducci\'on \(J=C_P\bmod3\) cumple
\[
J^2=2I_6=-I_6
\quad\text{en }\operatorname{End}_{\mathbb F_3}(\mathbb F_3^6).
\]
El cuarto de torsi\'on determina entonces la fase c\'iclica
\(1,\ii,-1,-\ii\). La aplicaci\'on del teorema chino del resto conserva
ambos datos residuales y produce \(Q_{12}=Q_3Q_4\). La traza
\(\Tr(N^{-s}Q)\) se eval\'ua en la etapa terminal; ni la serie de Dirichlet
ni su expresi\'on decimal intervienen en la selecci\'on del car\'acter.

La construcci\'on se explicita mediante las tablas
\begin{equation}
\chi_{-3}(n)=
\begin{cases}
0,&3\mid n,\\
1,&n\equiv1\pmod3,\\
-1,&n\equiv2\pmod3,
\end{cases}
\qquad
\chi_{-4}(n)=
\begin{cases}
0,&2\mid n,\\
1,&n\equiv1\pmod4,\\
-1,&n\equiv3\pmod4.
\end{cases}
\label{eq:characters34-tables}
\end{equation}
En \(\ell^2(\N)\),
\[
Q_3e_n=\chi_{-3}(n)e_n,
\qquad Q_4e_n=\chi_{-4}(n)e_n.
\]
Como son diagonales, conmutan; el teorema chino del resto da
\begin{equation}
(Q_3Q_4)e_n=\chi_{-3}(n)\chi_{-4}(n)e_n
=\chi_{12}(n)e_n.
\label{eq:q12-product-proof}
\end{equation}
El producto conserva simult\'aneamente el residuo tri\'adico y el cuarto de
torsi\'on. Por consiguiente, el car\'acter m\'odulo doce queda determinado
antes de efectuar cualquier evaluaci\'on de sus valores \(L\).

\begin{remark}[Control negativo]
La cadena queda invalidada si \(Q_3\) no distingue las dos clases orientadas,
si \(J^2\ne-I\), si \(Q_3Q_4\) no tiene la tabla CRT declarada o si el
operador se selecciona a partir del valor de la constante.
\end{remark}

\section{Estructura compleja interna y constante de Catalan}

La estructura compleja interna satisface

\begin{equation}
J^2=-I.
\label{eq:J-square}
\end{equation}

La igualdad \(J^2=-I\) define una representaci\'on
\(\rho:C_4\to\operatorname{Aut}_{\mathbb F_3}(\mathbb F_3^6)\),
\(\rho(j)=J\). Para realizarla espectralmente se toma el car\'acter unitario
\(\upsilon:C_4\to U(1)\), \(\upsilon(j)=\ii\), y el mapa de fase
\(n\mapsto j^n\). Sobre \(\ell^2(\N)\), la representaci\'on diagonal inducida
es
\[
U_4e_n=\upsilon(j^n)e_n=\ii^ne_n.
\]
Su parte orientada es
\[
Q_4=\frac{U_4-U_4^*}{2\ii},\qquad
Q_4e_n=\sin\!\left(\frac{\pi n}{2}\right)e_n
=\chi_{-4}(n)e_n.
\]
Este es el levantamiento expl\'icito entre el generador finito de orden cuatro y el operador
residual; no se identifica \(\mathbb F_9\) con \(\mathbb C\).

Sobre enteros impares, el mismo car\'acter puede escribirse

\[
\chi_J(n)=\ii^{n-1}=\chi_{-4}(n).
\]

Equivalentemente,

\begin{equation}
Q_4=\frac{U_4-U_4^*}{2\ii},
\qquad Q_4e_n=\chi_{-4}(n)e_n.
\label{eq:Q4}
\end{equation}

La constante de Catalan aparece como

\begin{equation}
G=\Tr(N^{-2}Q_4)=L(2,\chi_{-4})
=\sum_{n\ge0}\frac{(-1)^n}{(2n+1)^2}.
\label{eq:catalan}
\end{equation}

El levantamiento unitario del cuarto de torsi\'on genera exactamente el
car\'acter cuyo momento cuadr\'atico es la constante de Catalan.

\begin{proof}[Prueba operatoria]
El operador \(N^{-2}Q_4\) es traza--clase porque
\(\sum_{n\ge1}|\chi_{-4}(n)|n^{-2}<\infty\). Su traza es, por tanto, la
suma de los elementos diagonales. La tabla \cref{eq:characters34-tables}
anula los pares y alterna los impares:
\[
\Tr(N^{-2}Q_4)
=1-\frac1{3^2}+\frac1{5^2}-\frac1{7^2}+\cdots.
\]
La identidad con \(L(2,\chi_{-4})\) es la definici\'on de la serie de
Dirichlet; la igualdad con \(\operatorname{Im}\operatorname{Li}_2(\ii)\)
se obtiene tomando la parte imaginaria de
\(\sum_{n\ge1}\ii^n/n^2\).
\end{proof}

El valor terminal es
\[
G=0.9159655941772190150546035149\ldots.
\]
La estructura determinante es la identidad \(J^2=-I\), que induce
exactamente el signo alternante de las clases impares; la expansi\'on decimal
es su evaluaci\'on terminal.

\section{Composici\'on residual tri\'adico--cuaternaria}

El funcional residual m\'odulo tres proporciona \(\chi_{-3}\). Su producto
CRT con \(\chi_{-4}\) define

\begin{equation}
\chi_{12}=\chi_{-3}\chi_{-4}.
\label{eq:chi12}
\end{equation}

Sobre las unidades m\'odulo doce,

\[
\chi_{12}(1)=1,\quad
\chi_{12}(5)=-1,\quad
\chi_{12}(7)=-1,\quad
\chi_{12}(11)=1,
\]

y se anula fuera de ellas. Dos evaluaciones exactas son

\begin{align}
L(1,\chi_{12})&=\frac{\log(2+\sqrt3)}{\sqrt3},
\label{eq:L1chi12}\\
L(2,\chi_{12})&=\frac{\pi^2}{6\sqrt3}.
\label{eq:L2chi12}
\end{align}

La primera contiene la unidad hiperb\'olica

\begin{equation}
\sqrt3L(1,\chi_{12})=\arcosh2=\log(2+\sqrt3).
\label{eq:pell2}
\end{equation}

Junto con \(\log(3+2\sqrt2)=2\log(1+\sqrt2)\) y
\(\log(30+\sqrt{899})\), esta identidad determina una familia de escalas
de Pell cuyos miembros conservan sus respectivas procedencias algebraicas.

Las evaluaciones terminales certificadas son

\[
L(1,\chi_{12})=0.7603459963009463475310942549\ldots,
\]

\[
L(2,\chi_{12})=0.9497031262940093952634984917\ldots.
\]

\begin{proof}[Evaluaciones dodecaf\'asicas]
Agrupando por clases residuales,
\begin{align*}
L(1,\chi_{12})
&=\sum_{n\ge0}\left(
\frac1{12n+1}-\frac1{12n+5}
-\frac1{12n+7}+\frac1{12n+11}\right)\\
&=\int_0^1
\frac{1-x^4-x^6+x^{10}}{1-x^{12}}\,dx
=\int_0^1\frac{1-x^4}{1+x^6}\,dx.
\end{align*}
Una descomposici\'on real en los factores cuadr\'aticos de \(1+x^6\)
produce
\(\log(2+\sqrt3)/\sqrt3\). Para el segundo momento, si
\(\psi_1\) es la trigamma,
\begin{align*}
L(2,\chi_{12})
&=\frac1{144}\bigl[
\psi_1(1/12)-\psi_1(5/12)
-\psi_1(7/12)+\psi_1(11/12)\bigr]\\
&=\frac{\pi^2}{144}\left(
\csc^2\frac\pi{12}-\csc^2\frac{5\pi}{12}\right)
=\frac{\pi^2}{6\sqrt3},
\end{align*}
donde se ha usado
\(\psi_1(z)+\psi_1(1-z)=\pi^2\csc^2(\pi z)\). Las dos pruebas emplean
la tabla CRT y no un valor decimal objetivo.
\end{proof}

\section{Constantes de Euler y Stieltjes y transformada theta--Mellin}

Sea
\[
\vartheta(t)=\sum_{n\in\Z}e^{-\pi n^2t},\qquad t>0.
\]
Para \(\Re s>1\), el funcional theta--Mellin recupera la identidad convergente

\[
\pi^{-s/2}\Gamma\!\left(\frac{s}{2}\right)\zeta(s)
=\frac12\int_0^\infty
\bigl(\vartheta(t)-1\bigr)t^{s/2-1}\,dt.
\]

La continuaci\'on meromorfa se obtiene al partir la integral en \(t=1\), usar
\(\vartheta(t)=t^{-1/2}\vartheta(1/t)\) y separar los t\'erminos polares. Su expansi\'on en \(s=1\),

\begin{equation}
\zeta(s)=\frac1{s-1}+\sum_{j\ge0}
\frac{(-1)^j\gamma_j}{j!}(s-1)^j,
\label{eq:stieltjes}
\end{equation}

sit\'ua \(\gamma=\gamma_0\) y las constantes de Stieltjes en una misma
familia de momentos renormalizados. Esta organizaci\'on expresa una
procedencia anal\'itica com\'un, sin atribuir una interpretaci\'on f\'isica
id\'entica a todos sus niveles.

